\section{Results}
\label{sec:results}

\begin{table*}[t]
\centering
\caption{Register ablation on low-data CIFAR-100 (100 images/class, ViT-Tiny, 50 epochs). Mean $\pm$ population standard deviation over the three seeds \texttt{42}, \texttt{1337} and \texttt{3407}; best value per column in bold. Test metrics come from a single pass over the untouched 10{,}000-image test split. Validation loss is the value at the selected epoch and is therefore a selected minimum, not an unbiased estimate. The gap $\Delta\mathcal{L}=\mathcal{L}_{\mathrm{val}}-\mathcal{L}_{\mathrm{train}}$ is formed within each seed before averaging. Losses are label-smoothed ($\varepsilon=0.1$) cross-entropy, so they carry a constant positive floor and are comparable across arms but not against unsmoothed values.}
\label{tab:main_results}
\setlength{\tabcolsep}{5pt}
\footnotesize
\begin{tabular}{lccccc}
\toprule
\textbf{Configuration} & \textbf{Test Top-1 (\%)} & \textbf{Test Loss} & \textbf{Val Loss} & \textbf{Gap $\Delta\mathcal{L}$} & \textbf{$\bar{H}^{(12)}$ (bits)} \\
\midrule
Baseline ($K=0$) & \textbf{75.19 $\pm$ 0.24} & \textbf{1.6357 $\pm$ 0.0018} & 1.6829 $\pm$ 0.0609 & 0.8047 $\pm$ 0.0628 & 4.876 $\pm$ 0.115 \\
Registers ($K=1$) & 74.66 $\pm$ 0.42 & 1.6500 $\pm$ 0.0089 & 1.6545 $\pm$ 0.0380 & 0.7764 $\pm$ 0.0354 & 4.859 $\pm$ 0.101 \\
Registers ($K=4$) & 75.06 $\pm$ 0.33 & 1.6375 $\pm$ 0.0117 & \textbf{1.6492 $\pm$ 0.0536} & \textbf{0.7680 $\pm$ 0.0514} & \textbf{4.914 $\pm$ 0.031} \\
Registers ($K=8$) & 74.67 $\pm$ 0.12 & 1.6560 $\pm$ 0.0022 & 1.6602 $\pm$ 0.0268 & 0.7773 $\pm$ 0.0288 & 4.779 $\pm$ 0.091 \\
\bottomrule
\end{tabular}
\end{table*}


\subsection{No register configuration improves held-out accuracy}

Table~\ref{tab:main_results} reports every arm on both splits. The
register-free control attains the best held-out Top-1 accuracy,
$75.19 \pm 0.24\%$, and the best held-out loss, $1.6357 \pm 0.0018$. The nearest
register arm is $K=4$ at $75.06 \pm 0.33\%$; $K=1$ and $K=8$ trail by roughly
half a point.

Top-5 accuracy is the one place where a register arm leads, and we report it
rather than omit it: $K=4$ reaches $92.54 \pm 0.23\%$ against $92.43 \pm 0.14\%$
for the control. The lead does not hold up. Paired per seed the difference is
$+0.11$ pp with a standard deviation of $0.39$, positive in two seeds of three,
$t = 0.48$, $p = 0.681$ --- a spread three times the effect. The two arms that
lose Top-1 also lose Top-5, and lose it in every seed
($K=1$: $-0.18$ pp, 0/3; $K=8$: $-0.38$ pp, 0/3), so the ordering of the four
arms is unchanged by the metric.

Table~\ref{tab:paired_effects} makes the comparison per seed. Against the
control, $K=1$ loses accuracy in 3 of 3 seeds
($\bar{\delta} = -0.53$ pp, $p = 0.111$) and $K=8$ loses in 3 of 3
($\bar{\delta} = -0.51$ pp, $p = 0.096$). $K=4$ is the only arm that is not
consistently worse, and it is not better either: it wins on one seed of three,
with $\bar{\delta} = -0.13$ pp and $p = 0.466$. The honest summary is that
registers cost between zero and half a point of accuracy here, and gain nothing.

\begin{table}[t]
\centering
\caption{Paired per-seed effect of register tokens relative to the $K=0$ control. Each difference is taken within a seed and then reduced over the three seeds, so seed-to-seed variance cancels. ``Dir.'' counts how many of the three seeds moved in the direction the corresponding hypothesis predicts. $t$ is a paired $t$-statistic with $\mathrm{df}=2$ and $p$ is its two-sided value; at $n=3$ these are reported as effect-size evidence, and no correction for the fifteen comparisons is applied, so a single small $p$ should not be read as a confirmatory result.}
\label{tab:paired_effects}
\setlength{\tabcolsep}{3pt}
\footnotesize
\begin{tabular}{llccccc}
\toprule
\textbf{Metric} & \textbf{Arm} & \textbf{Mean $\Delta$} & \textbf{SD} & \textbf{Dir.} & \textbf{$t$} & \textbf{$p$} \\
\midrule
Test Top-1 (pp) & $K=1$ & -0.53 & 0.33 & 0/3 & -2.75 & 0.111 \\
 & $K=4$ & -0.13 & 0.25 & 1/3 & -0.89 & 0.466 \\
 & $K=8$ & -0.51 & 0.30 & 0/3 & -3.00 & 0.096 \\
\midrule
Test loss & $K=1$ & +0.0143 & 0.0119 & 0/3 & 2.08 & 0.173 \\
 & $K=4$ & +0.0018 & 0.0122 & 2/3 & 0.25 & 0.823 \\
 & $K=8$ & +0.0203 & 0.0042 & 0/3 & 8.37 & 0.014 \\
\midrule
Validation loss & $K=1$ & -0.0284 & 0.0416 & 2/3 & -1.18 & 0.359 \\
 & $K=4$ & -0.0337 & 0.0091 & 3/3 & -6.40 & 0.024 \\
 & $K=8$ & -0.0227 & 0.0418 & 2/3 & -0.94 & 0.447 \\
\midrule
Gen.\ gap $\Delta\mathcal{L}$ & $K=1$ & -0.0283 & 0.0494 & 2/3 & -0.99 & 0.426 \\
 & $K=4$ & -0.0367 & 0.0149 & 3/3 & -4.25 & 0.051 \\
 & $K=8$ & -0.0274 & 0.0417 & 2/3 & -1.14 & 0.373 \\
\midrule
$\bar{H}^{(12)}$ (bits) & $K=1$ & -0.0171 & 0.0434 & 2/3 & -0.68 & 0.566 \\
 & $K=4$ & +0.0375 & 0.1797 & 1/3 & 0.36 & 0.752 \\
 & $K=8$ & -0.0971 & 0.0456 & 0/3 & -3.69 & 0.066 \\
\bottomrule
\end{tabular}
\end{table}


\subsection{H1: the gap does narrow --- on the validation split only}

\begin{figure}[t]
\centering
\includegraphics[width=\columnwidth]{gen_gap_vs_registers.pdf}
\caption{Generalization gap and held-out test loss against register count. Error
bars are $\pm 1$ population standard deviation over the three seeds; the dotted
line marks the control arm. The two panels disagree: the quantity H1 predicts
should fall does fall, and the quantity that would make that fall meaningful
does not.}
\label{fig:gap}
\end{figure}

The generalization gap behaves exactly as H1 predicts. Every register arm sits
below the control, and $K=4$ is lowest at $0.7680 \pm 0.0514$ against
$0.8047 \pm 0.0628$. Paired per seed, $K=4$ narrows the gap in 3 of 3 seeds
($\bar{\delta} = -0.0367$, $t = -4.25$, $p = 0.051$) and lowers validation loss
in 3 of 3 seeds ($\bar{\delta} = -0.0337$, $t = -6.40$, $p = 0.024$). Taken
alone, this is the positive result the study set out to look for.

It does not survive contact with the test split. The same $K=4$ arm shows a
paired test-loss difference of $+0.0018$ with $t = 0.25$ and $p = 0.823$: not
merely insignificant, but centred on zero, with two of three seeds on either
side. Fig.~\ref{fig:gap} places the two panels side by side; the left panel
falls and the right panel does not follow it.

Two observations locate the cause. First, final training loss is effectively
identical across arms---$0.8782$, $0.8780$, $0.8812$ and $0.8829$ for
$K = 0, 1, 4, 8$, a spread of $0.005$ against a within-arm standard deviation of
$0.002$--$0.005$. A regularizer that reduced overfitting would be expected to
raise training loss while lowering validation loss; here the training term does
not move, so the entire gap reduction comes from the validation term. Second,
the validation split is small. Its between-seed dispersion in Top-1 accuracy is
$2.17$ pp in the control arm, against $0.24$ pp on the test split---a factor of
nine (Fig.~\ref{fig:accuracy}). A quantity that is both used for checkpoint
selection and measured on 1{,}000 images is a selected extremum of a noisy
estimator, and a consistent shift in it is not by itself evidence about
generalization.

\begin{figure}[t]
\centering
\includegraphics[width=\columnwidth]{accuracy_vs_registers.pdf}
\caption{Top-1 accuracy against register count on the two evaluation splits,
drawn on a shared vertical range. The validation panel is nearly uninformative
at this sample size; the test panel, on ten times the data, resolves the arms
and places the control first.}
\label{fig:accuracy}
\end{figure}

\subsection{H2: the $K=8$ arm degrades, consistent with capacity dilution}

$K=8$ is the one arm with a statistically consistent negative effect on
held-out data. Its paired test-loss difference is $+0.0203$ with the same sign
in 3 of 3 seeds, $t = 8.37$, $p = 0.014$---the smallest $p$-value in the study,
and the only one below $0.05$ on a test-split quantity. It also has the lowest
Top-5 accuracy of the four arms and the lowest final-block attention entropy,
the latter falling in 3 of 3 seeds ($\bar{\delta} = -0.0971$, $p = 0.066$). Its
Top-1 accuracy is not the lowest: $K=1$ is marginally worse
($74.66\%$ against $74.67\%$), a gap far inside the seed noise, so on Top-1 the
two arms are tied for last rather than ordered.

The accuracy response across $K$ is non-monotonic, as H2 predicts: it drops at
$K=1$, partially recovers at $K=4$, and drops again at $K=8$. We report the
non-monotonicity because it is what the hypothesis names, but with three seeds
per arm and a total range of $0.53$ pp we do not claim the shape of the curve is
resolved. What is resolved is the endpoint: eight registers on a $d = 192$
backbone are measurably worse than none.

\subsection{H3: no entropy-collapse effect}

\begin{figure}[t]
\centering
\includegraphics[width=\columnwidth]{entropy_vs_layer.pdf}
\caption{Layer-wise attention entropy. Left: absolute entropy per block, with
$\pm 1\sigma$ bands over seeds; the four arms are visually indistinguishable.
Right: the paired per-seed difference from the control, which is the only view
in which effects of this size are legible.}
\label{fig:entropy}
\end{figure}

Fig.~\ref{fig:entropy} shows the layer profile. Entropy is strongly
depth-dependent and non-monotonic---a minimum of $3.34$ bits at block 2, a
maximum of $6.27$ bits at block 5, and a decline through the final third---but
that profile is a property of the pretrained backbone, not of the treatment: the
four arms lie on top of one another. Averaged over all twelve blocks the paired
differences are $+0.020$, $+0.001$ and $-0.025$ bits for $K = 1, 4, 8$, none
significant. Registers do lower entropy in block 1 (by $0.19$ bits at $K=8$,
$0.12$ at $K=4$, $0.07$ at $K=1$) and raise it in blocks 9 and 11 at
$K \in \{1,4\}$ (largest: $+0.15$ bits at $K=1$, block 9), but no arm shows the deep-layer
entropy \emph{rise} relative to control that H3 requires. The one consistent
deep-layer effect runs the wrong way: $K=8$ lowers block-12 entropy in every
seed.

The patch-norm outlier rate agrees. Averaged over blocks it is
$1.227\%$, $1.246\%$, $1.250\%$ and $1.254\%$ for $K = 0, 1, 4, 8$. The paired
differences are positive for every register arm---the artifact rate rises
slightly rather than falling---though none approaches significance
($p = 0.24$ to $0.53$). We therefore find no evidence for the artifact-removal
mechanism in this regime, in either the entropy or the outlier readout.

\subsection{Optimisation behaviour}

\begin{figure}[t]
\centering
\includegraphics[width=\columnwidth]{loss_curves.pdf}
\caption{Seed-averaged loss trajectories, $\pm 1\sigma$ bands. Training curves
are superimposed across arms; validation curves are dominated by the noise of a
1{,}000-image split.}
\label{fig:curves}
\end{figure}

Fig.~\ref{fig:curves} confirms that the arms are not distinguishable during
optimisation either. The training trajectories coincide from the first epoch,
which is the expected consequence of adding at most 1{,}536 parameters to a
5.54\,M-parameter model. Validation trajectories are visibly noisy and cross
repeatedly. Best epochs cluster late---means of $49.0$, $47.3$, $49.0$ and
$44.7$ for $K = 0, 1, 4, 8$---so no arm was still improving materially at the
budget boundary, and the $K=8$ arm's earlier mean best epoch is consistent with
it having less to gain from the remaining schedule.
